\documentclass[a4paper,latin]{paper} 
\usepackage{babel}  
\usepackage[margin=2.5cm]{geometry}
\usepackage{graphicx}
\usepackage{lipsum}
\usepackage{xcolor}
\usepackage{booktabs}
\usepackage{hyperref}
\usepackage{listings}

\sectionfont{\large\sf\bfseries\color{black!70!blue}} 
\renewcommand\keywordname{Clavem verborum}
\title{ModBus Commander}
\subtitle{User Manual\\
\hfill\includegraphics[height=2cm]{honeywell-logo}
\vspace{-2cm}}
\author{Woody Wu} 
\institution{Honeywell HPS}

\begin{document} 
\maketitle 
\hrule 
\hrule\bigskip

\section{Synopis} 
\noindent\begin{lstlisting}[language=bash]
    Example: mbc -u 2
    Options:
      --help          Show help                                            [boolean]
      --version       Show version number                                  [boolean]
      --server, -s    server IP address                       [default: "localhost"]
      --port, -p      tcp port number of server                       [default: 502]
      --dev, -d       name of serial device
      --baudrate, -b  serial baudrate                                [default: 9600]
      --unit-id, -u   modbus server address (a.k.a, unit id)              [required]
      --file, -f      run script from file
      --gport, -g     tcp port to access web gui                     [default: 3000]
\end{lstlisting}

\section{Run From Command Line}

\subsection{Connect to TCP}
\noindent\begin{lstlisting}[language=bash]
  mbc -u 2 -s 192.168.0.100
\end{lstlisting}
By default, it connects to tcp port 502. The port can be changed by '-p'.

\subsection{Connect to RS485}
\noindent\begin{lstlisting}[language=bash]
  mbc -u 2 -d /dev/tty/USB0 -b 9600
\end{lstlisting}

\subsection{Examples Of Commands}
After the program started, it print a prompt. You can intput ModBus command to execute.
The following are some examples.

\subsubsection{Read holding registers}
\noindent\begin{lstlisting}[language=bash]
    fc3 0 100
\end{lstlisting}
This read 100 holding registers starting from address 0. The out put is like below:
\noindent\begin{lstlisting}[language=bash]
    0000: 6f53 5be8 18c7 2a15 6835 9415 ead8 32b9 0117 a660 c39f e30f c964 57be ac8b 8a89
    0010: 4245 cd57 6fb8 ea7b af1d 6faa 9de9 5cec 5258 510b 734d 2265 5122 c6ec 7965 6ef5
    0020: 9fd0 b05d 3abd 20bf 171c 819f b312 29e4 5350 5e97 571b b04a 8e36 f006 ddd2 8f7d
    0030: 2957 f67a f0d4 11bb 8208 a6e8 6e77 2015 f9ea a7d5 556d 3f9e 69a1 a41c ad75 0832
    0040: 055f 592b ecb4 f84b b07e 0d63 7b2c 2871 ad49 a440 495f 074d eec9 f250 29e1 ae0e
    0050: 876e 3f5d a750 b2a1 9059 73ef a4b9 324c c140 bdc8 a198 100f acdb 5c04 a6fd 6f64
    0060: 8448 8b81 12cf 1e00
\end{lstlisting}

\subsubsection{Write holding registers}
\noindent\begin{lstlisting}[language=bash]
    fc6 102 0x1234
\end{lstlisting}
This write value 0x1234 to holding register 102.

\subsubsection{Read coils}
\noindent\begin{lstlisting}[language=bash]
    fc1 16 32
\end{lstlisting}
This read 32 coils values starting from address 16. The output is,
\noindent\begin{lstlisting}[language=bash]
0000:           11110110 00010111 11011010 11100011
\end{lstlisting}

\subsection{Commands list}
Below lists all the commands:

\noindent\begin{tabular}{ll}
    fc1 & $\langle$start addr$\rangle$ $\langle$quantity of coils$\rangle$ \\
    fc2 & $\langle$start addr$\rangle$ $\langle$quatity of discrete inputs$\rangle$ \\
    fc3 & $\langle$start addr$\rangle$ $\langle$quantity of holding registers$\rangle$ \\
    fc4 & $\langle$start addr$\rangle$ $\langle$quantity of input registers$\rangle$ \\
    fc5 & $\langle$addr$\rangle$ $\langle$coil-value$\rangle$ \\
    fc6 & $\langle$addr$\rangle$ $\langle$holding-register-value$\rangle$ \\
    fc15 & $\langle$addr$\rangle$ $\langle$coil-value ...$\rangle$ \\
    fc16 & $\langle$addr$\rangle$ $\langle$holding-register-value ...$\rangle$ \\
    echo & string \\
    delay & msecs \\
\end{tabular}

\section{Run From WEB GUI}
You can also use its WEB GUI by opening http://localhost:3000.  It uses the same commands
set with the command line mode as described.  In WEB GUI mode, you can also watch program's
trace information.

The WEB GUI only tested in Google Chrome.

\includegraphics[width=\textwidth]{webgui.png}

\section{Command History}
In either command line mode or WEB GUI mode, all exected commands are saved and 
can be always brought back to execte again by pressing Up key on the keyboard.

\section{Scripting}
With the '-f' option in the command mode, the program can load a command file
and executes all the commands defined in the file one by one. Some special commands
such as \texttt{echo} and \texttt{delay} could be useful when put in a command
file.

\section{Installation}

You firstly need to install nodejs from \url{http://nodejs.org}, the current LTS version
is 10.16.0.

Then you can unarchive the distributed archive file to any folder and run the below
command inside the directory to get all the dependencies installed:
\noindent\begin{lstlisting}[language=bash]
  $ npm install
\end{lstlisting}

Now you can start the program inside the directory. On Linux:
\noindent\begin{lstlisting}[language=bash]
  $ ./bin/mbc [OPTIONS]
\end{lstlisting}

And on Windows:
\noindent\begin{lstlisting}[language=bash]
  c:\work\> src/mbc.js [OPTIONS]
\end{lstlisting}

\end {document}
